%=============================================================================!
% 10.0  CHAPTER OPENING                                                       !
%=============================================================================!
The earlier chapters followed a handful of atoms in empty space. To
represent a liquid or a crystal, we need a sample packed at the material's
density whose atoms have neighbours on every side. Repeating a simulation
cell supplies those neighbours, but also changes how we find separations
and sum the forces. We develop these operations first, including the
Ewald sum for charges, then use neighbour lists to avoid testing every
pair at every step.

With the forces defined, we can prepare the positions and velocities and
advance them with the velocity Verlet loop of Chapter~9. Comparisons with
an independent library check the implementation, and timings identify
where compiled code can reduce its cost. Chapter~11 starts from this loop
and asks how far the step can be lengthened by constraining stiff bonds
and evaluating slowly changing forces less often.
